\BourbakiExerciseGroup

\begin{enumerate}
\def\labelenumi{\arabic{enumi})}
\item
  Let A be a commutative ring. Show that, if every submodule of every free A-module is free, then every ideal of A is principal, and A is an integral domain (and consequently a principal ideal domain).
\item
  Let M be a module over a principal ideal domain A; suppose M is a direct sum of cyclic submodules. Show that every submodule N of M is a direct sum of cyclic submodules (if T is the torsion submodule of N, show first of all that N/T is a free module, using Th. 1, and deduce that T has a complement in N; then show that T is a direct sum of cyclic submodules, using Ex. 12, c) of VII, p.~57).
\item
  Let \((p_n)\) be a strictly increasing sequence of prime numbers such that \(p_n\) does not divide any of the numbers
\end{enumerate}

\[
a + b \left(1 + p _ {1} + p _ {1} p _ {2} + \dots + p _ {1} p _ {2} \dots p _ {n - 1}\right) \quad \text { for } \quad | a | \leqslant n \quad \text { and } \quad | b | \leqslant n.
\]

In the rational plane \(\mathbf{Q}^2\) , consider the \(\mathbf{Z}\) -module \(\mathbf{M}\) generated by the sequence \((x_n)\) defined as follows:

\[
x _ {0} = (1, 0), \quad x _ {1} = (0, 1), \quad x _ {n + 1} = p _ {n} ^ {- 1} \left(x _ {0} + x _ {n}\right) \quad \text { for } \quad n \geqslant 1.
\]

Show that M has rank 2, and that every rank 1 submodule of M is cyclic (notice that \(x_{0}\) and \(x_{n}\) form a basis for the submodule of M generated by the \(x_{i}\) with \(i \leqslant n\) ).

\begin{enumerate}
\def\labelenumi{\arabic{enumi})}
\setcounter{enumi}{3}
\tightlist
\item
  Let \(E\) be a torsion-free module over a principal ideal domain \(A\) .
\end{enumerate}

\begin{enumerate}
\def\labelenumi{\alph{enumi})}
\item
  Show that, if F is a pure submodule of E (VII, p.~55, Ex. 7), then E/F is torsion-free, and conversely.
\item
  Let F be a submodule of E generated by a maximal pseudofree family (VII, p.~55, Ex. 4); then F is a pure, free submodule of E. Show that, for all \(\dot{x} \in \mathrm{E} / \mathrm{F}\) , there exist a noninvertible element \(\alpha\) of A and an element \(\dot{y} \in \mathrm{E} / \mathrm{F}\) such that \(\dot{x} = \alpha \dot{y}\) , but that E/F is not necessarily a divisible module (cf.~Ex. 3).
\end{enumerate}

\begin{enumerate}
\def\labelenumi{\arabic{enumi})}
\setcounter{enumi}{4}
\tightlist
\item
  Let A be a principal ideal domain having a unique maximal ideal \(\mathbf{A}\pi\) (cf.~VII, p.~48, Ex. 4). Let E be a torsion-free A-module, and M a free submodule of E, such that E/M is torsion-free, divisible and of rank 1. Let \((e_{\alpha})\) be a basis of M, let \(\dot{u}\) be a nonzero element of E/M, and \(u\) an element of the coset \(\dot{u}\) ; for each integer \(n > 0\) , let \(u_{n}\) be an element of E such that \(u \equiv \pi^{n} u_{n} (\bmod M)\) . Put \(u = \pi^{n} u_{n} + \sum_{\alpha} \lambda_{n\alpha} e_{\alpha}\) .
\end{enumerate}

\begin{enumerate}
\def\labelenumi{\alph{enumi})}
\item
  Show that M is generated by M and the \(u_{n}\) , and that \(\lambda_{n\alpha} - \lambda_{n + 1,\alpha} \in \mathbf{A}\pi^n\) for all \(\alpha\) and for every integer \(n > 0\) (cf.~Ex. 4, a)).
\item
  Suppose in addition that A is complete with respect to the topology T defined by taking the ideals \(A\pi^{n}\) as a basis of open neighbourhoods of 0 in A (Gen.~Top., III, p.~5); put \(\lambda_{\alpha} = \lim_{n \to \infty} \lambda_{n\alpha}\) . Show that, if \(\mu\) and \(\mu_{\alpha}\) are such that \(\mu u - \sum_{\alpha} \mu_{\alpha} e_{\alpha}\) belongs to all the
\end{enumerate}

submodules \(\pi^n\mathrm{E}\) , then necessarily \(\mu_{\alpha} = \lambda_{\alpha}\mu\) for all \(\alpha\) .

\begin{enumerate}
\def\labelenumi{\arabic{enumi})}
\setcounter{enumi}{5}
\tightlist
\item
  Let \(\mathbf{A}\) be a principal ideal domain with a unique maximal ideal \(\mathbf{A}\pi\) , and complete with respect to the topology \(\mathcal{T}\) defined in Ex. 5.
\end{enumerate}

\begin{enumerate}
\def\labelenumi{\alph{enumi})}
\item
  Show that any torsion-free A-module E of finite rank, containing no nonzero divisible submodules, is free (consider a submodule M of E generated by a maximal pseudofree family; if \(M \neq E\) show, with the help of Ex. 4, b) and 5, b), that E contains a nonzero divisible submodule). Generalise this result to the case where E has countable rank (proceed by induction).
\item
  In the free module \(E = A^{(N)}\) (the direct sum of countably infinitely many copies of A), let \((a_{n})_{n \geq 0}\) be the canonical basis, and let \(e_{n} = a_{n-1} - \pi a_{n}\) for each integer \(n \geq 1\) . Show that the submodule M of E generated by the \(e_{n}\) is pure, and that the \(e_{n}\) form a maximal pseudofree family, but that E/M is a divisible module, and hence M does not admit a complement in E.
\end{enumerate}

\begin{enumerate}
\def\labelenumi{\arabic{enumi})}
\setcounter{enumi}{6}
\tightlist
\item
  Let \(p\) be a prime number and let \(\mathbf{A}\) denote the ring of rational numbers of the form \(r / s\) , where \(s\) is coprime to \(p\) (VII, p.~48, Ex. 4).
\end{enumerate}

Let \(u = (1,0)\) and \(v = (0,1)\) be the elements of the canonical basis of the vector space \(E = Q^2\) over \(Q\) . Recursively define a sequence of elements \(u_n\) of \(E\) by \(u_0 = u\) and \(u_{n}=pu_{n+1}+v\) for \(n\geqslant0\) . Let F be the A-submodule generated by v and the \(u_{n}\) ; show that F is not a free module, but contains no nonzero divisible submodule (cf.~Ex. 6, a)).

\begin{enumerate}
\def\labelenumi{\arabic{enumi})}
\setcounter{enumi}{7}
\tightlist
\item
  Let A be a principal ideal domain which is not a field.
\end{enumerate}

\begin{enumerate}
\def\labelenumi{\alph{enumi})}
\item
  Show that the product A-module \(A^{N}\) admits no countable system of generators. (If A is countable, consider \(\operatorname{Card}(A^{N})\) ; otherwise notice that, if for all \(x \in A\) we let \(v_{x}\) denote the element \((x^{n})_{n \geq 0}\) of \(A^{N}\) , then the \(v_{x}\) form a free system).
\item
  Let \(\pi\) be an irreducible element of A. Let S denote the submodule of \(\mathbf{A}^{\mathbf{N}}\) consisting of those sequences \(z = (z_{n})_{n\in \mathbb{N}}\) for which there exists a sequence \((k(n))_{n\in \mathbb{N}}\) of integers tending to \(+\infty\) (depending on \(z\) ), such that \(z_{n}\in \mathbf{A}\pi^{k(n)}\) for all \(n\) . Show that the \((\mathbf{A} / \pi)\) -vector space \(\mathrm{S} / \pi \mathrm{S}\) has a countable basis.
\item
  Deduce from a) and b) that \(A^{N}\) is not a free A-module. (Otherwise, S would be a free A-module with a countable basis; but \((x_{n})\mapsto(x_{n}\pi^{n})\) is an injective A-linear map from \(A^{N}\) into S.)
\end{enumerate}

\begin{enumerate}
\def\labelenumi{\arabic{enumi})}
\setcounter{enumi}{8}
\tightlist
\item
  Let \(\mathbf{A}\) be a principal ideal domain and let \(\mathbf{M}\) be the dual of the \(\mathbf{A}\) -module \(\mathbf{A}^{\mathbf{N}}\) . Let \(f: \mathbf{M} \to \mathbf{A}^{\mathbf{N}}\) be the transpose map of the canonical injection \(\mathbf{A}^{(\mathbf{N})} \to \mathbf{A}^{\mathbf{N}}\) (where \(\mathbf{A}^{\mathbf{N}}\) is identified with the dual of \(\mathbf{A}^{(\mathbf{N})}\) ).
\end{enumerate}

\begin{enumerate}
\def\labelenumi{\alph{enumi})}
\item
  Let \(\pi\) be an irreducible element of A and let \(\varphi: \mathbf{A}^{\mathbf{N}} / \mathbf{A}^{(\mathbf{N})} \to \mathbf{A}\) be a linear form; show that if \(u \in \mathbf{A}^{\mathbf{N}} / \mathbf{A}^{(\mathbf{N})}\) is the class of an element \((u_n) \in \mathbf{A}^{\mathbf{N}}\) such that \(u_n \in (\pi^n)\) for all \(n\) , then \(\varphi(u) = 0\) . If A contains two non-associate irreducible elements, deduce that \(f\) is injective. (Using the above and Bezout's identity, show that the dual of \(\mathbf{A}^{\mathbf{N}} / \mathbf{A}^{(\mathbf{N})}\) is zero.)
\item
  Give A the topology defined in Ex. 5 of VII, p.~60. Show that, if \(f(\mathbf{M})\) is not contained in \(\mathbf{A}^{(\mathbf{N})}\) , then A is complete. (Reduce the problem to showing that, if \(\varphi : \mathbf{A}^{\mathbf{N}} \to \mathbf{A}\) is a linear form such that \(\varphi(e_n) \neq 0\) for every element \(e_n\) of the canonical basis of \(\mathbf{A}^{(\mathbf{N})}\) , then \(\operatorname{Im}(\varphi)\) contains the sum of any series of elements of A which tends to 0 sufficiently quickly.)
\item
  If A contains two non-associate irreducible elements, then the canonical map from \(\mathbf{A}^{(N)}\) (resp. \(A^{N}\) ) into its double dual is bijective (use a) and b)).
\end{enumerate}

\begin{enumerate}
\def\labelenumi{\arabic{enumi})}
\setcounter{enumi}{9}
\item
  Let M be a nonzero Z-submodule of the ring of p-adic integers \(Z_{p}\) such that M is a pure Z-submodule of the Z-module \(Z_{p}\) (VII, p.~55, Ex. 7). Show that \(M + pZ_{p} = Z_{p}\) ; deduce that M is indecomposable (observe that \(Z_{p}/pZ_{p}\) is isomorphic to Z/pZ). Give an example of a decomposable Z-submodule of \(Z_{p}\) (observe that \(Z_{p}\) has the cardinality of the continuum).
\item
  Let A and B be two infinite sets of rational prime numbers, with no common elements, and not containing the number 5. Let \((e_i)_{1 \leqslant i \leqslant 4}\) be the canonical basis of the Q-vector space \(\mathbf{Q}^4\) .
\end{enumerate}

\begin{enumerate}
\def\labelenumi{\alph{enumi})}
\item
  Let M be the Z-submodule of \(Q^{4}\) generated by the elements \(e_{1}/m\) for \(m \in A\) ; let N be the Z-submodule of \(Q^{4}\) generated by the elements \(e_{2}/m\) with \(m \in A\) , \(e_{3}/n\) with \(n \in B\) , and \((e_{2} + e_{3})/5\) . Then \(M \cap N = 0\) . Put \(e_{1}' = 8e_{1} + 3e_{2}\) and \(e_{2}' = 5e_{1} + 2e_{2}\) . Let \(M'\) be the Z-submodule of \(Q^{4}\) generated by the elements \(e_{1}'/m\) with \(m \in A\) ; let \(N'\) be the Z-submodule of \(Q^{4}\) generated by the elements \(e_{2}'/m\) with \(m \in A\) , \(e_{4}/n\) with \(n \in B\) , and \((3e_{2}' + e_{3})/5\) . Show that \(M' \cap N' = 0\) , that \(M \oplus N = M' \oplus N'\) ; that M and \(M'\) are isomorphic and indecomposable; and that N and N' are indecomposable but not isomorphic (compare VII, p.~69, Ex. 23).
\item
  Let M be the Z-submodule of \(Q^{4}\) generated by the elements \(e_{1}/m\) with \(m \in A\) , \(e_{2}/n\) with \(n \in B\) , and \((e_{1} + e_{2})/5\) ; let N be the Z-submodule of \(Q^{4}\) generated by the elements \(e_{3}/m\) with \(m \in A\) , \(e_{4}/n\) with \(n \in B\) , and \((e_{3} + e_{4})/5\) . Put \(e_{1}' = 2e_{1} + e_{3}\) , \(e_{2}' = e_{2} + 3e_{4}\) , \(e_{3}' = 17e_{1} + 9e_{3}\) , and \(e_{4}' = e_{2} + 2e_{4}\) . Let \(M'\) be the Z-submodule of \(Q^{4}\) generated by the elements \(e_{1}'/m\) for \(m \in A\) , \(e_{2}'/n\) for \(n \in B\) , and \((e_{1}' + 2e_{2}')/5\) ; let \(N'\) be the Z-submodule of \(Q^{4}\) generated by the elements \(e_{3}'/m\) for \(m \in A\) , \(e_{4}'/n\) for \(n \in B\) , and \((e_{3}' + 2e_{4}')/5\) . Show that \(M \oplus N = M' \oplus N'\) , that M is isomorphic to N and \(M'\) is isomorphic to \(N'\) , but \(M'\) is not isomorphic to M.
\item
  Let M be the Z-submodule of \(Q^{4}\) generated by \(e_{1}/5^{n}\) ( \(n \geqslant 0\) ); let N be the submodule of \(Q^{4}\) generated by \(e_{2}/5^{n}\) , \(e_{3}/2^{n}\) , \(e_{4}/3^{n}\) ( \(n \geqslant 0\) ), \((e_{2} + e_{3})/3\) and \((e_{2} + e_{4})/2\) . Put \(e_{1}' = 3e_{1} - e_{2}\) and \(e_{2}' = 2e_{1} - e_{2}\) . Let \(M'\) be the Z-submodule of \(Q^{4}\) generated by the elements \(e_{1}'/5^{n}\) , \(e_{3}/2^{n}\) ( \(n \geqslant 0\) ) and \((e_{1}' - e_{3})/8\) ; let \(N'\) be the Z-submodule of \(Q^{4}\) generated by the elements \(e_{2}'/5^{n}\) , \(e_{4}/3^{n}\) ( \(n \geqslant 0\) ) and \((e_{2}' - e_{4})/2\) . Show that \(M \oplus N = M' \oplus N'\) , that \(M'\) and \(N'\) are indecomposable of rank 2, and that M has rank 1.
\end{enumerate}
% semantic-fragments: legacy-input-seam
\relax{}
